\documentclass[11pt,a4paper]{article}

% ---------- Encoding, language and typography ----------
\usepackage[utf8]{inputenc}
\usepackage[T1]{fontenc}
\usepackage{lmodern}
\usepackage[english]{babel}
\usepackage{microtype}

\usepackage[margin=2.5cm,headheight=15pt]{geometry}
\usepackage{xcolor}
\usepackage{titlesec}
\usepackage{enumitem}
\usepackage{array}
\usepackage{booktabs}
\usepackage{tabularx}
\usepackage{float}
\usepackage{graphicx}
\usepackage{listings}
\usepackage{mdframed}
\usepackage{pifont}
\usepackage{parskip}
\usepackage{fancyhdr}

% ---------- Colors ----------
\definecolor{accent}{RGB}{150,72,38}
\definecolor{darkgray}{RGB}{90,90,90}
\definecolor{lightgray}{RGB}{245,245,245}
\definecolor{codebg}{RGB}{30,30,30}
\definecolor{codetext}{RGB}{220,220,220}
\definecolor{codekw}{RGB}{240,160,90}
\definecolor{codecmt}{RGB}{125,165,120}
\definecolor{okgreen}{RGB}{35,120,55}
\definecolor{warnred}{RGB}{179,38,30}

\definecolor{stTodo}{RGB}{115,115,115}
\definecolor{stInProgress}{RGB}{40,110,180}
\definecolor{stReview}{RGB}{200,130,20}
\definecolor{stMerged}{RGB}{120,60,160}
\definecolor{stDeployed}{RGB}{30,140,90}
\definecolor{stTest}{RGB}{205,95,35}
\definecolor{stDone}{RGB}{35,110,55}

% ---------- Headings and page style ----------
\titleformat{\section}
  {\Large\bfseries\color{accent}}{\thesection}{0.8em}{}
\titlespacing*{\section}{0pt}{2.0ex plus .4ex}{1.0ex}
\titleformat{\subsection}
  {\large\bfseries\color{accent!80!black}}{\thesubsection}{0.7em}{}

\pagestyle{fancy}
\fancyhf{}
\fancyhead[L]{\footnotesize\sffamily\color{darkgray}Ticket Manager}
\fancyhead[R]{\footnotesize\sffamily\color{darkgray}\nouppercase{\leftmark}}
\fancyfoot[C]{\footnotesize\sffamily\color{darkgray}---~\thepage~---}
\renewcommand{\headrulewidth}{0.8pt}
\renewcommand{\headrule}{{\color{accent}\hrule width\headwidth height \headrulewidth}}

\setlist{itemsep=2pt,parsep=0pt,topsep=4pt}
\emergencystretch=1em
\raggedbottom

% ---------- Helper commands ----------
\newcommand{\ok}{\textcolor{okgreen}{\ding{51}}}
\newcommand{\no}{\textcolor{darkgray}{\ding{55}}}
\newcommand{\risk}{\textcolor{warnred}{\ding{55}}}
\newcommand{\statusbadge}[2]{%
  \colorbox{#1}{\makebox[3.0cm][c]{\color{white}\sffamily\bfseries\small #2}}}
\newcommand{\file}[1]{\texttt{\detokenize{#1}}}

% ---------- Code listing style ----------
\lstdefinestyle{sql}{
  language=SQL,
  backgroundcolor=\color{codebg},
  basicstyle=\ttfamily\footnotesize\color{codetext},
  keywordstyle=\color{codekw}\bfseries,
  commentstyle=\color{codecmt}\itshape,
  stringstyle=\color{codetext},
  numbers=left,
  numberstyle=\tiny\color{darkgray},
  numbersep=8pt,
  frame=none,
  breaklines=true,
  showstringspaces=false,
  tabsize=2,
  xleftmargin=2.6em,
  framexleftmargin=1.6em,
}
\lstdefinestyle{html}{
  language=HTML,
  backgroundcolor=\color{codebg},
  basicstyle=\ttfamily\footnotesize\color{codetext},
  keywordstyle=\color{codekw},
  commentstyle=\color{codecmt}\itshape,
  stringstyle=\color{codetext},
  numbers=left,
  numberstyle=\tiny\color{darkgray},
  numbersep=8pt,
  frame=none,
  breaklines=true,
  breakindent=1.5em,
  showstringspaces=false,
  tabsize=2,
  xleftmargin=2.6em,
  framexleftmargin=1.6em,
}
\lstset{style=sql}

% ---------- PDF metadata ----------
\usepackage[unicode,colorlinks=true,linkcolor=accent,urlcolor=accent,citecolor=darkgray]{hyperref}
\hypersetup{
  pdftitle={Ticket Manager - Full System Documentation},
  pdfauthor={Ekin},
  pdfsubject={Architecture, Security and Project Report},
  pdfkeywords={ticket manager, kanban, hono, sqlite, drizzle, better-auth, self-hosted},
}

\begin{document}

% ============================================================
% COVER
% ============================================================
\begin{titlepage}
\centering
\vspace*{2cm}
{\color{accent}\rule{0.6\textwidth}{1pt}}
\vspace{0.8cm}

{\Huge \bfseries Ticket Manager\par}

\vspace{0.4cm}
{\Large \color{darkgray} Full System Documentation\\ Architecture, Security and Project Report\par}

\vspace{0.8cm}
{\color{accent}\rule{0.6\textwidth}{1pt}}
\vspace{1.2cm}

{\large \itshape A zero-cost, self-hosted internal Jira / Linear alternative for\\
PlannedLost development tracking and multi-project management\par}

\vspace{2cm}
{\normalsize
\begin{tabular}{rl}
\textbf{Date:}      & \today \\
\textbf{Version:}   & v1.0 \\
\textbf{Author:}    & Ekin \\
\textbf{Status:}    & Implemented and running (development stage) \\
\end{tabular}\par}
\vfill
\end{titlepage}

\pagenumbering{roman}
\tableofcontents
\newpage
\pagenumbering{arabic}

% ============================================================
\section{Executive Summary}

Ticket Manager is a self-hosted, zero-cost work tracking system built for
\textbf{PlannedLost}, an independent game development effort. Its purpose is
twofold: (1) to track PlannedLost's own development with a Kanban workflow and
unique ticket IDs, and (2) to manage any number of additional projects from a
single installation. The long-term goal is to replicate the paid core of tools
such as \href{https://linear.app/}{Linear} on our own server, without license
fees and without handing data to a third party.

This document is a complete, ground-up survey of the implemented system: every
technology choice and its justification, every subsystem and where it lives in
the code, the security model, the features that were deliberately left out and
why, and a prioritized list of weaknesses and recommended next steps --- from
both an engineering and a management perspective.

\subsection{Current State at a Glance}

\begin{table}[H]
\centering
\small
\begin{tabularx}{\textwidth}{p{4.4cm}X}
\toprule
\textbf{Area} & \textbf{Status}\\
\midrule
Kanban board (7 columns, drag \& drop, inline editing) & \ok~implemented \\
Ticket IDs (\texttt{PREFIX-n}), types (Task/Bug), priorities & \ok~implemented \\
Rich-text descriptions, comments, per-ticket activity feed & \ok~implemented \\
File attachments (image/video, up to 100\,MB) & \ok~implemented \\
Cookie-session auth (Better Auth), invitation-only sign-up & \ok~implemented \\
Global roles + per-project roles, enforced server-side & \ok~implemented \\
Admin panel (users \& projects, bidirectional access management) & \ok~implemented \\
Global search (normalized IDs, description text, typo-tolerant suggestions) & \ok~implemented \\
E-mail transport abstraction (Resend in production, file outbox in dev) & \ok~implemented \\
Scheduled jobs: daily digest, weekly report, weekly security scan (OSV) & \ok~implemented \\
Automated test suite & \risk~missing \\
Version control (git) for this working directory & \risk~missing \\
CI/CD pipeline, production deployment configuration & \risk~missing \\
Password reset flow, API rate limiting, pagination & \risk~missing \\
\bottomrule
\end{tabularx}
\caption{Implementation status summary. Red crosses mark known gaps discussed
in Section~\ref{sec:gaps}.}
\label{tab:status}
\end{table}

% ============================================================
\section{System Components and Technology Choices}

An earlier revision of this document (v0.5) described a Cloudflare-centric
design: Pages Functions, D1, Wrangler and a react-admin management panel.
During implementation the architecture pivoted to a \textbf{self-hosted
Node.js server with a local SQLite file} and a \textbf{hand-written vanilla
JavaScript front-end}. The pivot happened for practical reasons: the
application is an internal tool for a single team, a local database file is
easier to inspect, back up and reason about than a remote D1 instance, and the
Cloudflare free tier's read quotas added complexity without adding value at
this scale. The pivot also removed an entire build step: the React admin
bundle (and with it react-admin, React, Vite --- 153 packages) was deleted and
replaced by server-rendered static pages plus vanilla ES modules.

\begin{table}[H]
\centering
\small
\begin{tabularx}{\textwidth}{p{2.6cm}p{4.6cm}X}
\toprule
\textbf{Layer} & \textbf{Technology} & \textbf{Rationale}\\
\midrule
Runtime      & Node.js 24 + \file{tsx} (dev)  & Single-language stack; \file{tsx watch} gives instant restart during development.\\
\addlinespace
HTTP server  & Hono 4.9 via \file{@hono/node-server} & Small, fast, Web-standard Request/Response API; the same routing code could later be redeployed to Cloudflare Workers if ever desired.\\
\addlinespace
Database     & SQLite (better-sqlite3 12) + Drizzle ORM 0.45 & Zero-administration relational store in \file{data/ticket-manager.db}; WAL journaling; synchronous driver is the fastest option for a single-node tool. Drizzle provides typed schema and queries.\\
\addlinespace
Auth         & Better Auth 1.3 & Cookie sessions, e-mail/password, Drizzle adapter; the \file{role} column is registered as a non-input additional field so clients can never set it.\\
\addlinespace
Front-end    & Vanilla ES modules + hand-written CSS design system & No build step at all; the server statically serves \file{client/}. A framework was judged unnecessary for $\sim$2{,}200 lines of UI code (see Section~\ref{sec:frontend}).\\
\addlinespace
E-mail       & Resend HTTP API behind a transport interface & Production delivery via Resend; development writes rendered mails to a local outbox directory instead of a catch-all SMTP server (Section~\ref{sec:mail}).\\
\addlinespace
Scheduling   & In-process \file{setInterval} tick (60\,s) & No external cron needed; jobs are idempotent per period (Section~\ref{sec:cron}).\\
\addlinespace
Vulnerability feed & OSV (\texttt{api.osv.dev}) batch query & Free, structured CVE data for the weekly dependency scan.\\
\bottomrule
\end{tabularx}
\caption{Technology choices and their justifications as implemented today.}
\label{tab:techstack}
\end{table}

\subsection{Why the Alternatives Were Not Chosen}

\begin{itemize}
  \item \textbf{react-admin (removed).} The original plan used react-admin for
  the super-admin panel. In practice the panel only ever needed two small
  resources (users, projects); react-admin brought Material UI theming that
  clashed with the application's own design system, plus a separate Vite build
  pipeline. It was replaced by $\sim$530 lines of vanilla JS in
  \file{client/js/pages/admin.js} that reuses the application's CSS components,
  and the dependency was removed entirely. Security was unaffected: the panel
  was always gated server-side (Section~\ref{sec:rbac}).
  \item \textbf{Cloudflare D1 / Pages Functions (postponed).} Still possible
  in the future --- Hono was chosen partly because its handlers are
  Web-standard and portable --- but local SQLite removes network latency from
  every query during development and makes backup a file copy.
  \item \textbf{React/Vue/Svelte for the main client (rejected).} The UI is a
  handful of pages with one interactive board. A framework would have added a
  build step, virtual-DOM reconciliation and $\sim$50\,MB of
  \file{node_modules} to solve problems this application does not have.
  State is kept in a single exported \file{state} object and pages re-render
  explicitly.
  \item \textbf{MailHog/Mailpit SMTP catcher (rejected for now).} The original
  plan used an SMTP catcher for development. The implemented transport
  abstraction is even simpler in a Node context: without
  \file{MAIL_TRANSPORT=resend} every mail is written as
  \file{.mail-out/<timestamp>-<to>-<subject>.html} and logged to the console.
  No extra container is required, and the rendered HTML can be opened
  directly in a browser.
\end{itemize}

\subsection{Code Base Overview}

\begin{table}[H]
\centering
\small
\begin{tabularx}{\textwidth}{p{5.2cm}p{1.6cm}X}
\toprule
\textbf{Location} & \textbf{Lines} & \textbf{Contents}\\
\midrule
\file{WEB/server/main.ts} & 1{,}004 & All HTTP endpoints (Hono), static file serving, upload handling.\\
\file{WEB/src/lib/*} & $\sim$700 & RBAC, invites, search, activity, notify, ticket IDs, mail transport, status model.\\
\file{WEB/src/db/*} & $\sim$250 & Drizzle schema (13 tables), SQLite bootstrap, migration runner.\\
\file{WEB/src/auth/server.ts} & 27 & Better Auth factory.\\
\file{WEB/src/cron/*} & $\sim$300 & Scheduler, idempotency guard, digest / weekly / security jobs.\\
\file{WEB/migrations/*.sql} & 5 files & Hand-written, ordered, auto-applied schema migrations.\\
\file{WEB/client/js/**} & $\sim$2{,}175 & Vanilla SPA: router, API wrapper, UI kit, 6 pages.\\
\file{WEB/client/css/app.css} & 535 & The entire design system (single light theme).\\
\file{WEB/templates/*.html} & 12 files & E-mail templates compiled into \file{generated-templates.ts}.\\
\file{WEB/seed/seed.ts} & $\sim$160 & Deterministic demo data (4 users, 2 projects, tasks, invitations).\\
\bottomrule
\end{tabularx}
\caption{Repository layout and approximate sizes.}
\label{tab:layout}
\end{table}

% ============================================================
\section{Application Architecture}

\subsection{Request Lifecycle}

Every request flows through the same pipeline in \file{server/main.ts}:

\begin{enumerate}
  \item \textbf{Static assets.} \file{/js/*} and \file{/css/*} are served with
  \file{Cache-Control: no-cache} in development so edits show up on refresh;
  other client assets are served statically from \file{client/}.
  \item \textbf{Auth routes.} \file{/api/auth/*} is delegated to Better Auth's
  handler. One exception is installed \emph{before} it: a middleware rejects
  every public \file{POST /api/auth/sign-up/email} unless it carries an
  internal secret header (Section~\ref{sec:auth}).
  \item \textbf{API routes.} Each handler is wrapped in \file{guard()}, which
  catches \file{ApiError} and converts it into a JSON error response with the
  right status code; unexpected errors become a generic 500 without leaking
  internals. Every protected handler starts by resolving the session
  (\file{getSessionUser}) and asserting a role
  (\file{requireSuperAdmin}/\file{requirePanelAccess}/\file{requireProjectRole}).
  \item \textbf{SPA fallback.} Any non-API GET serves
  \file{client/index.html}; the client router renders the correct page for
  paths such as \file{/projects/prj_pln} or \file{/my-tickets}.
\end{enumerate}

\subsection{HTTP API Surface}

\begin{table}[H]
\centering
\small
\begin{tabularx}{\textwidth}{p{5.6cm}p{2.2cm}X}
\toprule
\textbf{Endpoint} & \textbf{Method(s)} & \textbf{Access}\\
\midrule
\file{/api/auth/*} & POST & Better Auth (sign-in, sign-out, session). Public sign-up disabled.\\
\file{/api/projects} & GET, POST & Any authenticated (list is role-filtered) / create: ADMIN+\\
\file{/api/projects/:id} & GET, PATCH, DELETE & Project role / project admin / super admin\\
\file{/api/my-tickets} & GET & Any authenticated (assigned tickets across accessible projects)\\
\file{/api/search?q=} & GET & Any authenticated (results role-filtered; user search ADMIN+)\\
\file{/api/projects/:id/tasks} & GET, POST & Project role\\
\file{/api/projects/:id/tasks/:taskId} & PATCH, DELETE & Assignee or creator or project admin / delete: project admin\\
\file{/api/projects/:id/tasks/:taskId/events} & GET & Project role (activity feed)\\
\file{/api/projects/:id/tasks/:taskId/comments} & GET, POST, DELETE & Project role\\
\file{/api/projects/:id/media} & POST & Project role (multipart upload)\\
\file{/api/projects/:id/invitations} & GET, POST & Project role (read) / project admin (create)\\
\file{/api/invitations/accept} & POST & Public with valid token (project or account invite)\\
\file{/api/admin/users}, \file{/api/admin/users/:id} & GET, PATCH, DELETE & Super admin (user list, role change, delete, detail + memberships)\\
\file{/api/admin/users/:id/projects} & POST, DELETE & Super admin (membership upsert/remove from the user side)\\
\file{/api/admin/projects/:id/members} & GET, POST, DELETE & Super admin (same memberships from the project side)\\
\file{/api/admin/invites} & POST & Super admin (standalone account invitation)\\
\file{/api/dev/cron/:job} & POST & Super admin (manual job trigger for development)\\
\bottomrule
\end{tabularx}
\caption{HTTP API surface. Every write path re-validates authorization
server-side; the client-side checks are purely cosmetic.}
\label{tab:api}
\end{table}

% ============================================================
\section{Data Model and Migrations}
\label{sec:datamodel}

\subsection{Schema}

Better Auth owns the identity tables (\file{user}, \file{session},
\file{account}, \file{verification}); the application tables reference them.
The full application schema follows.

\begin{lstlisting}[language=SQL]
-- Projects
CREATE TABLE projects (
    id                      TEXT PRIMARY KEY,
    name                    TEXT NOT NULL,
    prefix                  TEXT NOT NULL UNIQUE,     -- e.g. WCE
    current_ticket_sequence INTEGER NOT NULL DEFAULT 0,
    created_at              INTEGER NOT NULL
);

-- Tickets
CREATE TABLE tasks (
    id          TEXT PRIMARY KEY,
    project_id  TEXT NOT NULL REFERENCES projects(id) ON DELETE CASCADE,
    ticket_id   TEXT NOT NULL UNIQUE,                     -- e.g. WCE-12
    title       TEXT NOT NULL,
    description TEXT,                                     -- sanitized rich text (HTML)
    status      TEXT NOT NULL DEFAULT 'TODO',
    type        TEXT NOT NULL DEFAULT 'TASK' CHECK (...), -- TASK | BUG
    priority    TEXT NOT NULL DEFAULT 'MEDIUM',           -- LOW..URGENT
    assignee_id TEXT REFERENCES user(id) ON DELETE SET NULL,
    created_by  TEXT REFERENCES user(id) ON DELETE SET NULL,
    position    REAL NOT NULL DEFAULT 0,                  -- intra-column order
    created_at  INTEGER NOT NULL,
    updated_at  INTEGER NOT NULL
);
CREATE INDEX idx_tasks_project_status ON tasks(project_id, status);

-- Project memberships and in-project roles
CREATE TABLE project_members (
    project_id TEXT NOT NULL REFERENCES projects(id) ON DELETE CASCADE,
    user_id    TEXT NOT NULL REFERENCES user(id)    ON DELETE CASCADE,
    role       TEXT NOT NULL DEFAULT 'MEMBER',            -- ADMIN | MEMBER
    PRIMARY KEY (project_id, user_id)
);

-- Project invitations (only the token HASH is stored)
CREATE TABLE invitations (
    id           TEXT PRIMARY KEY,
    project_id   TEXT NOT NULL REFERENCES projects(id) ON DELETE CASCADE,
    email        TEXT NOT NULL,
    token_hash   TEXT NOT NULL UNIQUE,
    invited_role TEXT NOT NULL DEFAULT 'MEMBER',          -- ADMIN | MEMBER
    expires_at   INTEGER NOT NULL,
    accepted_at  INTEGER,
    created_by   TEXT NOT NULL REFERENCES user(id) ON DELETE CASCADE,
    created_at   INTEGER NOT NULL
);

-- Standalone account invitations (no project attached)   [migration 0005]
CREATE TABLE account_invites (
    id          TEXT PRIMARY KEY,
    email       TEXT NOT NULL,
    token_hash  TEXT NOT NULL UNIQUE,
    expires_at  INTEGER NOT NULL,
    accepted_at INTEGER,
    created_by  TEXT NOT NULL REFERENCES user(id) ON DELETE CASCADE,
    created_at  INTEGER NOT NULL
);

-- Field-level ticket history (source for the activity feed and digests)
CREATE TABLE activity_log (
    id         TEXT PRIMARY KEY,
    task_id    TEXT NOT NULL REFERENCES tasks(id) ON DELETE CASCADE,
    actor_id   TEXT REFERENCES user(id) ON DELETE SET NULL,
    event_type TEXT NOT NULL,   -- CREATED | STATUS_CHANGED | TITLE_CHANGED |
                               -- DESCRIPTION_CHANGED | ASSIGNEE_CHANGED |
                               -- TYPE_CHANGED | PRIORITY_CHANGED
    old_status TEXT, new_status TEXT,
    old_value  TEXT, new_value  TEXT,
    created_at INTEGER NOT NULL
);
CREATE INDEX idx_activity_created ON activity_log(created_at);

-- Ticket comments
CREATE TABLE task_comments (
    id        TEXT PRIMARY KEY,
    task_id   TEXT NOT NULL REFERENCES tasks(id) ON DELETE CASCADE,
    author_id TEXT REFERENCES user(id) ON DELETE SET NULL,
    body      TEXT NOT NULL,
    created_at INTEGER NOT NULL
);

-- Cron idempotency: one row per (job, period)
CREATE TABLE job_runs (
    job_name   TEXT NOT NULL,
    period_key TEXT NOT NULL,
    created_at INTEGER NOT NULL,
    PRIMARY KEY (job_name, period_key)
);

-- Weekly security scan de-duplication
CREATE TABLE security_findings (
    id           TEXT PRIMARY KEY,
    cve          TEXT NOT NULL,
    package_name TEXT NOT NULL,
    severity     TEXT NOT NULL,
    fixed_version TEXT,
    reported_at  INTEGER NOT NULL,
    resolved_at  INTEGER,
    UNIQUE (cve, package_name)
);
\end{lstlisting}

Design notes:
\begin{itemize}
  \item \textbf{Deletion semantics are explicit everywhere.} Deleting a
  project cascades to tasks, memberships and invitations; deleting a user
  cascades sessions, accounts and memberships while tickets fall back to
  ``Unassigned'' (\file{ON DELETE SET NULL}). This is what makes the admin
  panel's delete actions safe with a single \file{DELETE} statement.
  \item \textbf{\file{position} is a \file{REAL}.} After a drag \& drop the
  new position is the midpoint between the neighbours, so re-ordering never
  renumbers a whole column (Section~\ref{sec:kanban}).
  \item \textbf{\file{ticket_id} is \texttt{UNIQUE}} --- the schema-level
  backstop behind the atomic counter (Section~\ref{sec:ticketid}).
\end{itemize}

\subsection{Migration System}

Migrations are plain SQL files in \file{WEB/migrations}, applied in
lexicographic order at every server start by \file{src/db/local.ts}; applied
names are recorded in a \file{_migrations} table. Five migrations exist
(\file{0001_init} \dots \file{0005_account_invites}). The approach is
deliberately primitive: no migration tool dependency, the SQL that will run is
exactly the SQL a human reviewed, and a fresh checkout reaches the current
schema by just starting the server. The trade-off is that there is no
down-migration; destructive changes require the documented
\file{npm run db:reset} + re-seed path, which is acceptable for an internal
tool at this stage but should be revisited before real data accumulates.

% ============================================================
\section{Authentication and Account Lifecycle}
\label{sec:auth}

\subsection{Sessions}

Better Auth is configured with e-mail + password and a Drizzle adapter over
the same SQLite database. Sessions are cookie-based (\file{credentials:
same-origin} on every client fetch). The \file{role} field is declared with
\texttt{input: false}, so even a forged sign-up request body cannot influence
it --- every account starts as \texttt{USER}.

\begin{mdframed}[linecolor=accent,linewidth=0.8pt,
                 backgroundcolor=lightgray,skipabove=10pt,skipbelow=10pt]
\small
\textbf{Public sign-up is disabled.} Better Auth would normally expose
\file{POST /api/auth/sign-up/email} to everyone. A Hono middleware in front of
the auth handler rejects it with \texttt{403 "Sign-up is by invitation only"}
unless the request carries the internal secret header. The invitation
acceptance flow calls Better Auth's \file{signUpEmail} \emph{in-process}
(server-side function call, not HTTP) and passes that header explicitly, so
the two invitation flows keep working while the open endpoint is closed.
\end{mdframed}

\subsection{Invitation Flows}

There are two invitation kinds, sharing the same security mechanics
(\file{invite.ts} / \file{account-invite.ts}):

\begin{enumerate}
  \item \textbf{Project invitation} (created by project admins or super
  admins, \file{POST /api/projects/:id/invitations}): targets an e-mail
  address, carries an in-project role (\texttt{ADMIN}/\texttt{MEMBER}), and on
  acceptance the account is created (if needed) \emph{and} the membership row
  is inserted.
  \item \textbf{Account invitation} (created only by super admins,
  \file{POST /api/admin/invites}): creates the account with no project
  access; memberships are granted afterwards from the admin panel. This was
  added so that accounts can exist before any project assignment.
\end{enumerate}

Both share the same token discipline:
\begin{itemize}
  \item The raw token is 32 bytes from \file{crypto.getRandomValues}
  (CSPRNG), base64url-encoded; it exists \emph{only} inside the e-mail link.
  \item The database stores the \textbf{SHA-256 hash} with a \texttt{UNIQUE}
  constraint. A leaked database cannot be replayed as invitation links.
  \item Tokens expire after 48 hours and are single-use
  (\file{accepted_at} is stamped on acceptance).
  \item Duplicate pending invitations and already-existing accounts are
  rejected with \texttt{409}.
\end{itemize}

The acceptance endpoint (\file{POST /api/invitations/accept}) transparently
handles both kinds: it first validates the token as a project invitation and,
if that fails, as an account invitation. The invited person always sets
\emph{their own} password on the acceptance page --- no administrator ever
handles a user's password.

\subsection{Known Gaps in the Account Lifecycle}

\begin{itemize}
  \item \textbf{No password reset.} A forgotten password currently requires a
  super admin to intervene at the database level. This is the most important
  missing account-management feature.
  \item \textbf{No e-mail verification} independent of invitations; the
  \file{email_verified} column exists but is not part of any flow.
  \item \textbf{No rate limiting} beyond Better Auth's built-in defaults on
  auth endpoints; application API endpoints are unthrottled.
\end{itemize}

% ============================================================
\section{Role-Based Access Control}
\label{sec:rbac}

\subsection{Roles}

\begin{enumerate}
  \item \textbf{Super Admin} (\file{SUPER_ADMIN}): full control --- user
  list/roles/deletion, account invitations, project creation/deletion,
  membership management from either side, manual cron triggers.
  \item \textbf{Admin} (\file{ADMIN}, global): may create projects, invite
  people to projects they can access and view the user list read-only in the
  admin panel. \emph{Note:} in the current model a global Admin implicitly has
  admin rights on \emph{every} project (\file{requireProjectRole} short-circuits
  membership checks for non-\texttt{USER} roles). Tightening this to explicit
  per-project membership is a documented open decision
  (Section~\ref{sec:gaps}).
  \item \textbf{User} (\file{USER}): sees and acts only inside projects where
  a \file{project_members} row exists.
\end{enumerate}

Inside a project, members additionally carry \file{ADMIN} or \file{MEMBER}:

\begin{table}[H]
\centering
\small
\begin{tabular}{lccc}
\toprule
\textbf{Capability} & \textbf{Super Admin} & \textbf{Project Admin} & \textbf{Member}\\
\midrule
View board and tickets                     & \ok & \ok & \ok\\
Create tickets                             & \ok & \ok & \ok\\
Edit tickets they created / are assigned to& \ok & \ok & \ok\\
Edit \emph{any} ticket                     & \ok & \ok & \no\\
Reassign tickets                           & \ok & \ok & \no\\
Change ticket type or priority             & \ok & \ok & \no\\
Delete tickets                             & \ok & \ok & \no\\
Invite to project / view invitation list   & \ok & \ok & \no\\
Rename project, change prefix              & \ok & \ok & \no\\
Delete project                             & \ok & \no & \no\\
Create projects (global)                   & \ok & \ok & \no\\
Manage global roles / delete users         & \ok & \no & \no\\
\bottomrule
\end{tabular}
\caption{Effective permission matrix. ``Project Admin'' includes global
Admins and Super Admins, who short-circuit membership checks.}
\label{tab:rbac}
\end{table}

\subsection{Enforcement}

Enforcement is exclusively server-side, in three layers:
\begin{enumerate}
  \item \file{getSessionUser} rejects anonymous requests (\texttt{401}).
  \item \file{requireSuperAdmin} / \file{requirePanelAccess} gate the admin
  API (\texttt{403}).
  \item \file{requireProjectRole} resolves the effective in-project role and
  rejects non-members (\texttt{403}); ticket-level rules (assignee/creator
  edits, admin-only fields) are checked inside the task PATCH handler.
\end{enumerate}
The client hides controls the current user cannot use, but this is cosmetic:
every capability denial is a server decision, so IDOR-style attempts through
the API fail regardless of the UI.

% ============================================================
\section{Kanban Workflow and Ordering}
\label{sec:kanban}

\subsection{Columns}

\begin{table}[H]
\centering
\small
\begin{tabularx}{\textwidth}{clX}
\toprule
\textbf{Order} & \textbf{Status} & \textbf{Meaning}\\
\midrule
1 & \statusbadge{stTodo}{TO DO}             & Backlog; accepted work waiting to be picked up.\\
2 & \statusbadge{stInProgress}{IN PROGRESS} & Actively being worked on.\\
3 & \statusbadge{stReview}{UNDER REVIEW}    & Code complete, awaiting review. Instant e-mail.\\
4 & \statusbadge{stMerged}{MERGED}          & Integrated into the main branch.\\
5 & \statusbadge{stDeployed}{DEPLOYED}      & Live. Instant e-mail.\\
6 & \statusbadge{stTest}{TEST}              & Verification on the live system. Instant e-mail.\\
7 & \statusbadge{stDone}{DONE}              & Closed. Feeds the daily digest.\\
\bottomrule
\end{tabularx}
\caption{Kanban columns.}
\label{tab:kanban}
\end{table}

\subsection{Transitions and Ordering}

\texttt{canTransition()} currently allows \emph{any} status to any status.
This is a documented, deliberate decision: the board is drag \& drop based and
free movement makes accidental drags trivially reversible. The server still
validates status values and rejects same-status no-ops that are not reorders.
If stricter flow control is wanted later, the single function in
\file{src/lib/status.ts} is the only place to change.

Ordering uses fractional indexing: when a ticket is dropped before another
one, its new \file{position} is the midpoint between the target and its
predecessor (\file{computePosition}). Dropping at the end takes
\file{max(position) + 1}. Reordering therefore costs one \file{UPDATE} and
never renumbers a column. The practical limit is float precision after
$\sim$50 consecutive inserts at the same spot, which is far beyond observed
usage; a periodic renumber could be added if it ever matters.

% ============================================================
\section{Ticket Identity Generation}
\label{sec:ticketid}

Each project owns a unique prefix (\file{WCE}, \file{PLN}, \dots); tickets are
\file{<PREFIX>-<n>}. The counter is incremented atomically in a single
statement so two concurrent creations can never receive the same number:

\begin{lstlisting}[language=SQL]
UPDATE projects
SET current_ticket_sequence = current_ticket_sequence + 1
WHERE id = ?
RETURNING current_ticket_sequence;
\end{lstlisting}

SQLite serializes writes, and the \texttt{UNIQUE} constraint on
\file{tasks.ticket_id} is the final backstop. Note that the sequence counts
\emph{creations}, not existing tickets: deleted tickets are not re-used and
the counter is not decremented. The UI therefore shows the real
\texttt{COUNT(*)} of tickets per project (added to \file{/api/projects}) next
to the sequence-derived ID space.

% ============================================================
\section{Ticket Content: Rich Text, Comments, History, Media}

\begin{itemize}
  \item \textbf{Rich-text descriptions} are entered in a contenteditable
  editor and sanitized on the client with an allowlist sanitizer
  (\file{sanitizeDesc}): only formatting tags survive, \file{<script>} and
  event handlers are stripped, links get \file{rel="noopener noreferrer"},
  and embedded \file{<img>}/\file{<video>} sources must point at
  \file{/media/uploads/*}. Legacy plain-text descriptions are auto-upgraded.
  The sanitizer is client-side only; server-side sanitization is a listed gap
  (Section~\ref{sec:gaps}).
  \item \textbf{Comments} are stored per ticket (\file{task_comments}) with
  \file{@mention} support in the composer (member autocomplete popup).
  \item \textbf{Activity feed.} Every mutation writes a row to
  \file{activity_log} (seven event types). The feed endpoint resolves
  assignee IDs into display names and is rendered as the right-hand column of
  the ticket modal. The daily digest and weekly report query the same table.
  \item \textbf{Media.} \file{POST /api/projects/:id/media} accepts multipart
  uploads; the MIME type is checked against a whitelist (png, jpg, gif, webm,
  mp4, mov), the size against 100\,MB, and the file is stored under a random
  UUID name in \file{data/uploads}, served back at \file{/media/uploads/*}.
  There is no orphan cleanup: deleting a ticket does not remove its embedded
  files (Section~\ref{sec:gaps}).
\end{itemize}

% ============================================================
\section{Global Search}

The header search (\file{GET /api/search}) was built to tolerate the way
people actually type:

\begin{enumerate}
  \item \textbf{Fast path (SQL).} The query is normalized (lowercase,
  non-alphanumerics stripped), and tickets are matched with
  \file{LIKE} against the normalized \file{ticket_id} expression
  (\file{replace(lower(ticket_id), '-', '')}), so \texttt{pln2} finds
  \texttt{PLN-2}. Titles \emph{and} rich-text descriptions are matched with a
  plain \file{LIKE} on the raw query. Single-character queries are restricted
  to ticket IDs and project prefixes to avoid matching half the database.
  \item \textbf{Fuzzy fallback.} If the fast path finds nothing, up to 500
  tickets from the caller's accessible projects are scored in JavaScript with
  a bounded Levenshtein distance (long queries require distance $\leq$1,
  short ones $\leq$2; a sentinel value keeps ``too far'' results out).
  Results are returned with a \texttt{suggested: true} flag and the UI labels
  them ``Did you mean''.
  \item \textbf{Projects and users.} Projects match on name/prefix within the
  accessible set; user matching (name/e-mail) is only executed for
  ADMIN+ callers --- a plain member's search response never even contains the
  user group.
\end{enumerate}

All results are filtered by the caller's accessible project set, so search
cannot leak tickets from projects the user cannot open. Verified examples:
\texttt{pln2} $\rightarrow$ PLN-2; \texttt{kanbn} $\rightarrow$ PLN-2
(suggested); \texttt{compose} $\rightarrow$ PLN-9 (description match).

% ============================================================
\section{Admin Panel}

The admin area lives inside the main application at \file{/admin/users} and
\file{/admin/projects}, rendered by the same vanilla SPA shell (the sidebar
stays visible; the ``Admin Panel'' entry is shown only to super admins).

\begin{itemize}
  \item \textbf{Users tab.} A read-only table; clicking a row opens a
  ``Manage user'' dialog with the user's global role selector (self-edit and
  self-delete are blocked server-side as lockout protection) and the full
  project access matrix: per project, a checkbox and a project-role selector.
  Saving sends only the diff (membership upserts/removals + optional role
  change). A danger zone deletes the account; foreign keys cascade sessions
  and memberships and unassign tickets.
  \item \textbf{Projects tab.} Lists all projects with real ticket counts;
  ``+ New project'' creates one (name + unique 2--5 character prefix,
  validated server-side). Clicking a row opens the mirrored ``Project
  access'' dialog: all users with checkboxes and project roles, with global
  Admins/Super Admins pinned as ``global access''. A delete action requires
  confirmation and cascades everything.
  \item \textbf{Invite user.} Super admins can invite an e-mail address
  without attaching a project (Section~\ref{sec:auth}); in development the
  dialog shows the generated link for copying since no SMTP server runs.
\end{itemize}

Both dialogs edit the same \file{project_members} table, so the two views
cannot drift apart by construction.

% ============================================================
\section{E-mail Infrastructure}
\label{sec:mail}

\subsection{Transport Abstraction}

All outgoing mail goes through a single \file{MailTransport} interface
(\file{src/lib/mail/index.ts}):

\begin{itemize}
  \item \textbf{\file{resend}} (production): posts to Resend's REST API; the
  API key is read from the environment, never committed.
  \item \textbf{\file{file}} (development, the default): writes the rendered
  HTML and plain-text parts into \file{WEB/.mail-out/} and logs the path.
  This replaces the originally planned MailHog container with zero
  infrastructure: the development ``inbox'' is a directory of openable files,
  and it is physically impossible to e-mail a real user by accident.
\end{itemize}

Selection is one environment variable (\file{MAIL_TRANSPORT}); switching to
production delivery changes no business code.

\subsection{Templates}

Twelve HTML templates live in \file{WEB/templates/} and are compiled into
\file{src/lib/mail/generated-templates.ts} by \file{scripts/gen-templates.ts}
(\file{npm run gen:templates}, run automatically before dev/seed). They share
a common skeleton: brand header in the accent colour (\#964826), table-based
responsive layout, inline styles for client compatibility, status-coloured
badges, and \texttt{\{\{placeholder\}\}} substitution with an
\file{htmlToText} plain-text fallback for every message.

\begin{table}[H]
\centering
\small
\begin{tabularx}{\textwidth}{p{5.2cm}X}
\toprule
\textbf{Template} & \textbf{Used by}\\
\midrule
\file{invite.html} & Project invitation e-mail.\\
\file{account-invite.html} & Standalone account invitation.\\
\file{status-<status>.html} (7) & Instant notifications for \texttt{UNDER REVIEW}, \texttt{DEPLOYED}, \texttt{TEST}.\\
\file{digest-daily.html} & Per-project daily summary (closed + in-progress).\\
\file{digest-weekly.html} & Weekly progress report to admins.\\
\file{security-report.html} & Weekly OSV scan findings.\\
\bottomrule
\end{tabularx}
\caption{E-mail template inventory.}
\label{tab:templates}
\end{table}

% ============================================================
\section{Scheduled Jobs}
\label{sec:cron}

\subsection{Scheduler}

\file{src/cron/scheduler.ts} runs a 60-second in-process tick. The current
wall-clock time is converted to \file{Europe/Istanbul} with
\file{Intl.DateTimeFormat} (no external timezone library), so ``midnight''
means Turkish midnight regardless of the host's timezone or DST:

\begin{table}[H]
\centering
\small
\begin{tabularx}{\textwidth}{p{3.4cm}p{3.4cm}X}
\toprule
\textbf{Job} & \textbf{Schedule (TR time)} & \textbf{Behaviour}\\
\midrule
Daily digest & every day 00:00 & Per project: tickets that reached \texttt{DONE} that day (from \file{activity_log}) plus all \texttt{IN PROGRESS} tickets; sent to project members.\\
\addlinespace
Weekly report & Sunday 22:00 & Progress roll-up sent to global Admins and Super Admins.\\
\addlinespace
Security scan & Sunday 21:00 & OSV batch query over the package list in \file{OSV_PACKAGES}; HIGH/CRITICAL findings are de-duplicated against \file{security_findings} and mailed to super admins with affected package, CVE, severity and fixed version.\\
\bottomrule
\end{tabularx}
\caption{Scheduled jobs. All are also triggerable by a super admin via
\file{POST /api/dev/cron/:job} for testing.}
\label{tab:cron}
\end{table}

\subsection{Idempotency}

Every run claims its period with \file{runOnce()}: an
\file{INSERT ... ON CONFLICT DO NOTHING} into \file{job_runs} keyed by
\file{(job_name, period_key)}. If the process restarts mid-day, the digest
is not sent twice; if the machine is off at the scheduled minute, the next
tick with a fresh period key runs it once. The security scan additionally
suppresses re-reporting of open findings via the
\file{security_findings} unique index.

% ============================================================
\section{Front-End Architecture}
\label{sec:frontend}

\subsection{SPA Shell}

\file{client/js/main.js} implements a history-API router with six routes
(\file{/login}, \file{/accept-invite}, \file{/projects},
\file{/projects/:id}, \file{/projects/:id/members}, \file{/my-tickets},
\file{/admin/*}). After authentication every page renders inside a persistent
\emph{shell}: a full-height project sidebar (My Tickets, project list, and an
Admin Panel entry for super admins) plus the page area. The header carries the
brand, the global search (centered, Ctrl+K (Cmd+K on macOS) focusable, debounced dropdown
with grouped results and ``did you mean'' suggestions) and the account box.

\subsection{Design System}

\file{client/css/app.css} ($\sim$535 lines) defines the whole visual language
with CSS variables (accent \#964826, neutral grays), reusable components
(buttons, chips, pills, tables, modals, toasts, tabs) and the Kanban styles.
Because the admin panel was migrated into this system, the entire application
has one coherent look with zero component-framework overhead. The theme is
light-only by design for now; the variables would make a dark theme a
mechanical change.

\subsection{State and Data}

State is deliberately simple: a module-level \file{state} object (current
user, project cache) exported from \file{main.js}, and explicit re-renders.
\file{api.js} is a thin \file{fetch} wrapper (same-origin credentials, JSON
errors surfaced as \file{Error.status}). There is no client-side cache
invalidation problem because there is almost no cache; the trade-off is more
refetching, which is irrelevant at LAN latency.

% ============================================================
\section{Security Assessment}
\label{sec:security}

\subsection{Implemented Hardening}

\begin{itemize}
  \item \textbf{Server-side authorization on every endpoint} (three-layer
  checks, Section~\ref{sec:rbac}); object IDs are always scoped by project or
  user, preventing IDOR.
  \item \textbf{Invitation tokens}: CSPRNG, 48-hour expiry, single-use,
  SHA-256-at-rest (Section~\ref{sec:auth}).
  \item \textbf{Public sign-up disabled} at the router level with an internal
  secret required to create accounts.
  \item \textbf{\file{role} cannot be set by clients} (Better Auth
  \texttt{input: false}).
  \item \textbf{Lockout protections}: a super admin can neither change their
  own role nor delete their own account.
  \item \textbf{Upload hardening}: MIME whitelist, size cap, random file
  names, no execution context (static file directory).
  \item \textbf{Client-side XSS defence} for ticket descriptions
  (allowlist sanitizer) and blanket \file{esc()} escaping of interpolated
  strings in the UI layer.
  \item \textbf{Weekly automated vulnerability scanning} against OSV with
  e-mailed HIGH/CRITICAL findings and de-duplication.
  \item \textbf{Secrets via environment}: \file{RESEND_API_KEY},
  \file{BETTER_AUTH_SECRET}, \file{MAIL_FROM}, \file{APP_URL},
  \file{OSV_PACKAGES}; a warning is printed when the auth secret falls back
  to a development default.
\end{itemize}

\subsection{Open Weaknesses}

\begin{table}[H]
\centering
\small
\begin{tabularx}{\textwidth}{p{4.6cm}p{1.8cm}X}
\toprule
\textbf{Weakness} & \textbf{Severity} & \textbf{Detail / recommended fix}\\
\midrule
No automated tests & High & Nothing prevents a regression in the atomic counter, RBAC or invite flows. Priority: integration tests over the API with an in-memory SQLite database (Vitest + Supertest-style harness).\\
\addlinespace
Not under version control & High & The working directory is not a git repository; a mistake or disk failure is unrecoverable. Priority: \file{git init} + private remote, immediately.\\
\addlinespace
No CI/CD & Medium & Lint, typecheck (\file{npm run check}), tests and builds should run on push.\\
\addlinespace
Client-side-only description sanitization & Medium & The server stores what it receives; a crafted API call could persist unsafe HTML. Fix: run the same allowlist sanitizer server-side before insert/update.\\
\addlinespace
No API rate limiting & Medium & Auth endpoints have Better Auth defaults; everything else is unthrottled. Fix: a small in-memory limiter per session/IP on write endpoints.\\
\addlinespace
No password reset & Medium & Manual database intervention required today. Fix: Better Auth's reset plugin or an invitation-style token flow.\\
\addlinespace
No pagination & Low & Task lists and the fuzzy pool (500) are bounded by convention, not by page parameters; fine for hundreds of tickets, not thousands.\\
\addlinespace
Orphaned media files & Low & Attachments of deleted tickets remain on disk. Fix: periodic reconciliation of \file{data/uploads} against referenced URLs.\\
\addlinespace
No HTTPS/deployment story & Medium & The server binds plain HTTP on :8788. Production needs a reverse proxy (Caddy/nginx) or TLS termination, a process manager, and a documented backup procedure for \file{data/}.\\
\addlinespace
Single light theme, no a11y audit & Low & Cosmetic/accessibility; CSS variables make theming cheap.\\
\bottomrule
\end{tabularx}
\caption{Known weaknesses with recommended remediations.}
\label{tab:weaknesses}
\end{table}

% ============================================================
\section{Deliberate Non-Decisions}
\label{sec:nondone}

The following were consciously \emph{not} built, with reasons recorded so
they are not mistaken for oversights:

\begin{itemize}
  \item \textbf{No front-end framework.} $\sim$2{,}200 lines of vanilla JS
  cover six pages; a framework's build step and reconciliation were judged
  costlier than explicit re-renders. Revisit only if the UI grows
  substantially (e.g. real-time multi-user boards).
  \item \textbf{No real-time collaboration.} Board updates require a refresh;
  there is no WebSocket/SSE layer. Acceptable for a single small team; the
  next meaningful upgrade would be SSE for board invalidation.
  \item \textbf{Free status transitions.} Any column to any column, so drags
  are undoable. A stricter machine can be introduced later in one function.
  \item \textbf{No per-project membership requirement for global Admins.}
  Global Admins implicitly administer every project. This keeps the seed team
  frictionless but is flagged as an open model decision (the admin UI already
  communicates it as ``global access'').
  \item \textbf{No Docker.} The app is one \file{npm run dev} command plus a
  SQLite file; containerization is deferred until a second machine deploys.
  \item \textbf{No i18n.} The UI is uniformly English; all strings are
  inline, so extraction would be required before localizing.
  \item \textbf{No SMTP catcher in development.} The file transport covers
  the same need with less infrastructure (Section~\ref{sec:mail}).
\end{itemize}

% ============================================================
\section{Gaps and Recommended Roadmap}
\label{sec:gaps}

\subsection{Engineering}

\begin{enumerate}
  \item \textbf{Put the project under git now} (including this document);
  without it, every other improvement is built on sand.
  \item \textbf{Integration test suite} for the API: RBAC matrix, invite
  lifecycle (expiry, reuse, duplicate), atomic ticket IDs under concurrency,
  search behaviour, cron idempotency. The synchronous SQLite driver makes an
  in-memory test database trivial.
  \item \textbf{CI pipeline}: typecheck, tests, \file{npm audit} on every
  push.
  \item \textbf{Server-side sanitization} of rich text before persistence.
  \item \textbf{Password reset flow} (token-based, mirroring the invite
  mechanics that already exist).
  \item \textbf{Rate limiting} on write endpoints and invitation creation.
  \item \textbf{Pagination parameters} on task/comment/event lists.
  \item \textbf{Media garbage collection} for unreferenced uploads.
  \item \textbf{Decide the global-Admin access model}: implicit
  all-project access (today) vs. explicit memberships with only
  \file{SUPER_ADMIN} exempt. The admin UI is already prepared for either.
  \item \textbf{Deployment recipe}: reverse proxy + TLS, process manager,
  documented backup/restore of \file{data/} (SQLite file + uploads).
\end{enumerate}

\subsection{Management}

\begin{itemize}
  \item \textbf{Bus factor is 1.} Everything is documented here, but only one
  person knows the codebase. The roadmap above (tests + CI + git) is also the
  onboarding path for a second contributor.
  \item \textbf{Data is the asset.} The SQLite file plus \file{data/uploads}
  is the entire state; a scheduled offline copy (even \file{rsync} to another
  disk) should exist from day one of real usage.
  \item \textbf{Scope discipline.} The pivot away from Cloudflare and
  react-admin both \emph{reduced} complexity; the same discipline should
  apply to future features --- the current feature set is complete for a
  small team, and the highest-value next investments are reliability
  (tests/CI/backups), not features.
  \item \textbf{Cost.} Running cost is effectively zero (self-hosted, Resend
  free tier); the only recurring obligation is the weekly security scan
  review.
\end{itemize}

% ============================================================
\section{Version History}

\begin{table}[H]
\centering
\small
\begin{tabularx}{\textwidth}{p{1.6cm}p{2.6cm}X}
\toprule
\textbf{Version} & \textbf{Date} & \textbf{Notes}\\
\midrule
v0.1--v0.5 & --- & Architecture drafts (see \file{DOC/main (1).tex}): Cloudflare D1/Pages design, react-admin panel, MailHog plan, e-mail template system, security scanning concept.\\
\addlinespace
v1.0 & \today & First documentation of the \emph{implemented} system: self-hosted Hono + SQLite pivot; vanilla SPA shell with persistent sidebar; admin panel migrated in-app (react-admin and Vite removed, 153 packages dropped); account invitations with public sign-up disabled; bidirectional user/project access management; global search with normalization, description matching and fuzzy suggestions; My Tickets cross-project view; real ticket counts; full security assessment and roadmap.\\
\bottomrule
\end{tabularx}
\caption{Document version history.}
\end{table}

\end{document}
